\section{模型评价与展望}
\subsection{模型优点}
\begin{enumerate}
\item 三问使用同一完整方案 $x=(\mathcal P,a,\boldsymbol\pi)$ 与字典序目标。子图划分、核心映射和核内顺序通过各问的执行映射共同作用于总体任务执行时间。第一问的子图级 Task、第二问的核心级 Task 以及第三问的 FIFO 缓存，都能在统一方案空间内比较。

\item 目标值由完整事件执行确定。片上 L1/UB 按缓冲实例检查瞬时容量，自动换出按固定规则插入。矩阵、向量及读写流水线各受占用约束，DDR 与 L2 带宽按在途请求共享。通信量和活跃峰值用于候选排序，下界用于筛选，完整事件执行确定最终周期。第三问的同方案配对固定子图划分、核心映射和核内顺序，使缓存对照具有明确的比较对象。

\item 多种候选来源共同改进一个已经验证的可行方案。结构移动提供并行度调整，驻留候选响应张量生命周期和跨核通信，缓存时序候选定位真实的未命中、填充与淘汰窗口。每次只接受完整评价可行且目标字典序改善的方案，保证已保留结果的可行性和单调性。全局下界给出相对当前周期仍可能降低的有效上界，可用于判定近界停止，并量化停止时的质量范围。
\end{enumerate}

\subsection{模型缺点}
\begin{enumerate}
\item 完整评价成本较高。通过结构预检、去重和下界筛选而进入完整评价的候选，需要重建任务、COPY、换入换出和全局事件序列。在共享带宽与 FIFO 状态耦合较强的图上，细小的顺序移动也可能引发大量事件变化。图规模、请求密度和评价成本因而限制可考察的候选数量。

\item 存储与带宽估计对事件时序的表达有限。按逻辑活跃集估计的峰值不能精确替代按缓冲实例计量的 L1/UB 占用。静态跨核字节也不包含读写请求的同时在途程度。第三问的读取可能发生在填充前或淘汰后，静态估计因而可能错排候选。

\item 来源调用记录中的改善次数受调用顺序和当前方案影响，难以单独量化各邻域的贡献。后续可固定初始方案与预算，逐项关闭邻域进行对照。六图容量与带宽研究也可扩展到全部算例，以考察参数响应的普遍性。
\end{enumerate}

\subsection{模型的推广与改进}
可进一步发展针对事件耦合的增量评价。候选只移动少量完整子图时，先确定受影响的依赖后继、同核 Pipe 队列和张量生命周期，再从最早变化事件向后更新。若能证明未波及部分的状态和就绪时刻不变，即可避免每次从头执行全图。增量结果仍应与完整事件语义等价，并在状态变化越过核心或带宽池时扩大重算范围。结合更紧的候选下界，可用节省的时间考察更多候选。

对缓存问题，可把轨迹中的读入发射、完成填充和淘汰构成显式窗口图，估计候选移动对关键请求所在窗口的影响。估计应同时考虑两带宽池、读入流水线和关键链。推广到允许选择缓存替换、预取或张量重算策略的场景时，可在现有 $x$ 之外增加相应决策变量，并重新证明可行性与下界。
